
\documentclass[11pt,a4paper]{letter}

\usepackage[margin=1in]{geometry}
\usepackage[T1]{fontenc}
\usepackage{lmodern}
\usepackage{microtype}
\usepackage{enumitem}
\usepackage[hidelinks]{hyperref}

\setlength{\parindent}{0pt}
\setlength{\parskip}{7pt}

\begin{document}

\begin{flushright}
9 November 2026
\end{flushright}

\textbf{To}\\
The Editor-in-Chief\\
\textit{Quantum Machine Intelligence}\\
Springer Nature

\vspace{0.3cm}

\textbf{Subject: Submission of Manuscript --- Hybrid Quantum--Classical
Pauli-Correlator Architecture for Image Reconstruction}

\vspace{0.3cm}

Dear Editor-in-Chief,

We are pleased to submit our manuscript entitled
\textbf{``Hybrid Quantum--Classical Pauli-Correlator Architecture
for Image Reconstruction''}
for consideration for publication in
\textit{Quantum Machine Intelligence}.

In this work, we introduce the
\textbf{Pauli-Correlator Autoencoder (PCAE)},
a hybrid quantum--classical framework combining matrix product
state (MPS) compression, variational quantum circuits, physically
motivated Hamiltonian regularization, and classical neural
networks for image reconstruction.

Our central objective is to investigate how quantum-mechanical
observables, particularly local Pauli expectation values and
two-qubit correlations, can serve as structured latent
representations for reconstructing classical visual information.

The principal contributions of our manuscript are:

\begin{enumerate}[leftmargin=*, itemsep=4pt]

\item \textbf{Hybrid quantum--classical architecture:}
We develop two six-qubit architectures, PCAE1D and PCAE2D,
employing MPS-based feature extraction and Pauli-correlator
representations with different circuit connectivity patterns.

\item \textbf{Image reconstruction across multiple datasets:}
We evaluate the framework on six image datasets, including
CIFAR-10, MNIST, Fashion-MNIST, and three medical imaging
datasets, demonstrating reconstruction of diverse image
structures through per-image optimization.

\item \textbf{Symmetry-aware quantum representations:}
We introduce and numerically verify a group-averaged observable
representation that is equivariant under the dihedral group
$D_4$, providing a mathematically controlled approach to
incorporating rotational and reflection symmetries.

\item \textbf{Experimental validation on IBM quantum hardware:}
We execute trained models on real IBM Quantum processors,
reconstructing images from experimentally measured quantum
observables without retraining the classical decoder.
The hardware experiments achieve reconstruction PSNR values
ranging from 25.90 to 31.52 dB in the initial five-image pilot.

\item \textbf{Comparative evaluation and physical interpretation:}
We perform classical baseline comparisons, component
ablations, noise analyses, and controlled tests of pairwise
observables to investigate the representational capabilities
and limitations of the architecture.

\end{enumerate}

We believe that this manuscript is particularly well suited
to \textit{Quantum Machine Intelligence}, as it addresses
the intersection of quantum information processing,
tensor-network representations, quantum machine learning,
and experimentally implementable hybrid computational
architectures.

Beyond reconstruction feasibility, our study examines the
structure and limitations of quantum-derived feature
representations, including their behavior under realistic
hardware noise.

Importantly, we do not claim a quantum computational advantage
over classical approaches. Instead, we present a systematic
investigation of how structured quantum observables can be
incorporated into classical learning pipelines, supported by
quantitative comparisons and real-hardware experiments.

We thank you for considering our manuscript for publication
in \textit{Quantum Machine Intelligence}.
We look forward to your evaluation.

\vspace{0.4cm}

Sincerely,

\textbf{Sparsho Chakraborty}\\
Department of Physics, School of Basic Sciences\\
Indian Institute of Technology Bhubaneswar\\
Odisha, India

\vspace{0.2cm}

\textit{On behalf of the authors:}\\
Sparsho Chakraborty and Siddhartha Patra

\vspace{0.2cm}

\textbf{Corresponding Author:}\\
Dr.\ Siddhartha Patra\\
Centre for Quantum Engineering, Research and Education\\
TCG CREST, Kolkata, India\\
\href{mailto:siddhartha.patra@tcgcrest.org}
{siddhartha.patra@tcgcrest.org}

\end{document}
